% =======================================================================
%  Related Papers Report
%  Two selected papers on perceptual HRTF evaluation metrics
%  -----------------------------------------------------------------------
%  Compile:  pdflatex related_papers_report
%            bibtex   related_papers_report
%            pdflatex related_papers_report
%            pdflatex related_papers_report
% =======================================================================
\documentclass[12pt,a4paper]{article}
\usepackage[utf8]{inputenc}
\usepackage[T1]{fontenc}
\usepackage{lmodern}
\usepackage[margin=2.5cm]{geometry}
\usepackage{hyperref}
\usepackage{booktabs}
\usepackage{enumitem}
\usepackage{parskip}
\usepackage{setspace}
\usepackage{amsmath}
\usepackage[numbers]{natbib}
\setstretch{1.5}

\begin{document}

% ===  TITLE  =============================================================
\pagenumbering{gobble}
\thispagestyle{empty}

\vspace*{2cm}

\begin{center}
{\LARGE \textbf{Related Work Report:\\
Two Selected Papers on Perceptual HRTF Evaluation Metrics}}
\end{center}

\vspace{1cm}

\begin{center}
Shanghao Zou\\[0.3cm]
{\small Prepared for the thesis\\
\textit{Objective Head-Related Transfer Function Evaluation Metrics\\
for More Consistent Perceptual Assessment}}
\end{center}

\vspace{0.5cm}

\begin{center}
\today
\end{center}

\newpage
\pagenumbering{arabic}

% ===  INTRODUCTION  ======================================================
\section{Introduction}

This report reviews two recent publications that are closely related to the
thesis topic of improving objective HRTF evaluation metrics.  Both papers
address the well-documented failure of Log-Spectral Distortion (LSD) as a
perceptual predictor and propose distinct solutions:

\begin{enumerate}[label=(\roman*)]
  \item \textbf{Yao et~al.\ (2024)} introduce a signal-processing pipeline
        called \emph{Perceptually Enhanced Spectral Distance} (PESD) that
        augments spectral comparison with auditory-model stages~\cite{yao2024pesd}.
  \item \textbf{Zhang et~al.\ (2025)} take a deep-learning approach,
        embedding HRTFs into a perception-informed latent space supervised by
        Metric Multidimensional Scaling (MMDS)~\cite{zhang2025perception}.
\end{enumerate}

For each paper the report covers the research motivation, datasets, models and
algorithms, key quantitative results, and relevance to the thesis.


% =========================================================================
\newpage
\section{Paper 1: Yao et~al.\ (2024) --- PESD}
% =========================================================================

\subsection{Bibliographic Information}

\begin{quote}
D.~Yao, J.~Zhao, Y.~Liang, Y.~Wang, J.~Gu, M.~Jia, H.~Lee, and J.~Li,
``Perceptually enhanced spectral distance metric for head-related transfer
function quality prediction,''
\textit{J.\ Acoust.\ Soc.\ Am.}, vol.~156, no.~6, pp.~4133--4152, Dec.~2024.
DOI: \href{https://doi.org/10.1121/10.0034632}{10.1121/10.0034632}.
\end{quote}

% ---  Motivation  --------------------------------------------------------
\subsection{Research Motivation}

LSD is the most widely used objective metric for evaluating HRTF quality,
yet multiple studies have shown its poor correlation with perceptual
judgements~\cite{doma2023, ananthabhotla2021}.  Yao et~al.\ identify three
specific shortcomings of LSD:

\begin{enumerate}
  \item \textbf{No perceptual thresholds.}  LSD penalises all spectral
        differences equally, even those below the auditory
        just-noticeable difference~(JND).
  \item \textbf{Uniform frequency weighting.}  The human cochlea analyses
        sound on a roughly logarithmic frequency scale; LSD treats every
        frequency bin identically.
  \item \textbf{No directional sensitivity.}  Perceptual importance of
        spectral cues varies across spatial directions, but LSD is
        direction-agnostic.
\end{enumerate}

% ---  Dataset  -----------------------------------------------------------
\subsection{Datasets}

\begin{table}[h]
\centering
\caption{Dataset used by Yao et~al.\ (2024).}
\label{tab:pesd_data}
\begin{tabular}{@{}ll@{}}
\toprule
\textbf{Property}        & \textbf{Detail} \\ \midrule
Primary database          & Institute of Acoustics, Chinese Academy of Sciences \\
Number of subjects        & 128 \\
Spatial positions         & 734 per subject \\
Listening-test directions & 32 \\
LSD distortion levels     & 6 (ranging from 1 to 6\,dB) \\
Test protocol             & MUSHRA-like \\
Perceptual attributes     & Spatial quality, timbral quality \\
\bottomrule
\end{tabular}
\end{table}

The listening tests applied controlled magnitude-spectrum distortions to
measured HRTFs at six LSD levels across 32~spatial directions, and
participants rated both spatial and timbral attributes using a MUSHRA-like
interface.

% ---  Algorithm  ---------------------------------------------------------
\subsection{Model / Algorithm --- The PESD Pipeline}

PESD replaces the single-step LSD calculation with a five-stage
auditory-model pipeline:

\begin{enumerate}
  \item \textbf{Spectral analysis.}
        The HRTF magnitude spectra are decomposed through perceptually
        motivated filter banks (analogous to cochlear frequency analysis),
        replacing the uniform FFT-bin treatment of LSD.

  \item \textbf{Threshold discrimination.}
        Spectral differences that fall below the auditory JND are zeroed
        out, preventing the metric from penalising imperceptible
        distortions.

  \item \textbf{Feature combination.}
        The surviving supra-threshold spectral differences are aggregated
        into a unified feature vector that emphasises perceptually salient
        frequency regions.

  \item \textbf{Binaural weighting.}
        Interaural cues (ITD and ILD) are incorporated via direction-dependent
        weights, reflecting the distinct roles of spatial and timbral
        attributes in HRTF perception.

  \item \textbf{Perceptual outcome estimation.}
        The weighted features are mapped to a scalar quality score that
        predicts the perceived similarity between a reference HRTF and a
        distorted version.
\end{enumerate}

% ---  Results  -----------------------------------------------------------
\subsection{Key Results}

\begin{itemize}
  \item PESD achieves \textbf{significantly higher Pearson correlation}
        with listening-test ratings than LSD, spectral distortion~(SD),
        and Mel-frequency cepstral distance~(MFCD).
  \item PESD exhibits \textbf{lower prediction error} (RMSE) across both
        spatial and timbral quality dimensions.
  \item The threshold-discrimination stage alone accounts for a
        substantial improvement, confirming that sub-threshold penalties
        are a major source of LSD's poor perceptual alignment.
  \item Results are consistent across the 32~evaluated directions,
        demonstrating the benefit of direction-dependent binaural weighting.
\end{itemize}

% ---  Relevance  ---------------------------------------------------------
\subsection{Relevance to the Thesis}

The thesis proposes a Composite Perceptual Metric (CPM) that integrates
LSD with ILD error, group delay error, and ITD error.  PESD pursues the
same goal---closing the gap between objective measurement and perceptual
quality---but through a fundamentally different strategy: rather than
combining multiple independent metrics, it re-engineers the spectral
distance computation itself with auditory-model stages.  PESD therefore
serves as both a \emph{benchmark} against which CPM can be compared and a
source of design principles (threshold gating, binaural weighting) that
could inform future iterations of CPM.


% =========================================================================
\newpage
\section{Paper 2: Zhang et~al.\ (2025) --- Perception-Informed Latent
         HRTF Representations}
% =========================================================================

\subsection{Bibliographic Information}

\begin{quote}
Y.~Zhang, A.~Francl, R.~Gao, P.~Calamia, Z.~Duan, and
I.~Ananthabhotla,
``Towards perception-informed latent HRTF representations,''
in \textit{Proc.\ IEEE WASPAA}, 2025 (Best Student Paper Award).
arXiv: \href{https://arxiv.org/abs/2507.02815}{2507.02815}.
\end{quote}

% ---  Motivation  --------------------------------------------------------
\subsection{Research Motivation}

Deep-learning methods for HRTF personalisation typically compress measured
HRTFs into a low-dimensional latent space via autoencoders trained to
minimise spectral reconstruction error (e.g., Spectral Difference
Error, SDE).  Zhang et~al.\ demonstrate that such latent spaces are
\emph{not} perceptually faithful: two HRTFs that are close in latent
distance can sound significantly different, and vice versa.  This
disconnect limits the usefulness of latent-space nearest-neighbour
selection for HRTF personalisation.

% ---  Dataset  -----------------------------------------------------------
\subsection{Datasets}

\begin{table}[h]
\centering
\caption{Datasets used by Zhang et~al.\ (2025).}
\label{tab:zhang_data}
\begin{tabular}{@{}lll@{}}
\toprule
                          & \textbf{SS2 (primary)} & \textbf{HUTUBS (generalisation)} \\ \midrule
Number of subjects        & 78                     & 96 \\
Training / test split     & 65 / 13                & --- (test only) \\
Purpose                   & Training \& evaluation & Cross-dataset generalisation \\
\bottomrule
\end{tabular}
\end{table}

The \textbf{Sound Sphere~2 (SS2)} database provides the primary training
and evaluation data, with 65~subjects for training and 13~for testing.
The \textbf{HUTUBS} database is used solely to evaluate cross-dataset
generalisation.

% ---  Architecture  ------------------------------------------------------
\subsection{Model / Algorithm}

\subsubsection{Baseline: Spectral Autoencoder}

A standard autoencoder compresses per-direction HRTF magnitude spectra
into a latent vector~$\mathbf{z}$ by minimising spectral reconstruction
loss (L2 / SDE).  The authors show that this baseline yields limited
Pearson correlation between latent distances and three objective
perceptual metrics.

\subsubsection{Proposed: INR-Based Perception-Informed Framework}

\begin{enumerate}
  \item \textbf{Implicit Neural Representation (INR).}
        Following the ``HRTF~Field'' paradigm, a multilayer perceptron (MLP)
        with 2048~hidden units models the HRTF magnitude as a continuous
        function of spatial direction, avoiding the discretisation artefacts
        of grid-based representations.

  \item \textbf{Perception-informed latent space via MMDS.}
        Instead of minimising spectral reconstruction error alone, the
        latent code~$\mathbf{z}$ is supervised by a
        \emph{Metric Multidimensional Scaling} (MMDS) loss that forces
        pairwise latent distances to match pairwise perceptual distances
        derived from objective metrics.

  \item \textbf{Perceptual distance targets.}
        The MMDS supervision uses distances computed from three objective
        perceptual metrics:
        \begin{itemize}
          \item \textbf{Coloration} --- Predicted Binaural Coloration (PBC),
          \item \textbf{Externalization} --- Auditory Externalization
                Perception (AEP),
          \item \textbf{Localization} --- Difference of Root Mean Square
                Error (DRMSE) in polar angles.
        \end{itemize}
\end{enumerate}

% ---  Results  -----------------------------------------------------------
\subsection{Key Results}

\begin{itemize}
  \item \textbf{Baseline (SDE-optimised) correlations:}
        Pearson~$r$ between latent distances and perceptual distances was
        limited---reported values of approximately 0.78, 0.73, and
        0.37 for coloration, localization, and externalization,
        respectively.
  \item \textbf{MMDS-supervised model:}
        Significantly higher Pearson correlations across all three
        perceptual dimensions, confirming that latent distances now
        better predict perceptual similarity.
  \item \textbf{HUTUBS generalisation:}
        Improvements transfer to the unseen HUTUBS dataset, indicating
        that the perception-informed structure is not dataset-specific.
  \item \textbf{HRTF selection:}
        When the latent space is used for nearest-neighbour HRTF
        personalisation, the perception-informed model selects HRTFs that
        are perceptually closer to the listener's own HRTF compared to the
        SDE-only baseline.
\end{itemize}

% ---  Relevance  ---------------------------------------------------------
\subsection{Relevance to the Thesis}

Zhang et~al.'s work intersects with the thesis in two ways:

\begin{enumerate}
  \item \textbf{Metric alignment.}
        Their finding that spectral reconstruction error alone is a poor
        proxy for perceptual quality directly mirrors the thesis's
        critique of LSD.  The three perceptual metrics they employ
        (PBC, AEP, DRMSE) offer an alternative or complementary
        evaluation vocabulary to the CPM components (LSD, ILD~error,
        group delay error, ITD~error).

  \item \textbf{Downstream application.}
        The MMDS-supervised latent space could serve as a downstream
        consumer of whatever composite metric the thesis develops:
        if CPM proves to be a more reliable perceptual distance than
        any single metric, it could replace or augment the PBC / AEP /
        DRMSE targets used by Zhang et~al., further improving
        personalisation.
\end{enumerate}


% =========================================================================
\newpage
\section{Comparative Summary}
% =========================================================================

\begin{table}[h]
\centering
\caption{Side-by-side comparison of the two reviewed papers.}
\label{tab:comparison}
\begin{tabular}{@{}p{3.5cm}p{5.5cm}p{5.5cm}@{}}
\toprule
& \textbf{Yao et~al.\ (PESD)} & \textbf{Zhang et~al.\ (MMDS)} \\
\midrule
\textbf{Core idea}
  & Re-engineer spectral distance with auditory-model stages
  & Supervise latent HRTF space with perceptual distances \\[6pt]
\textbf{Approach}
  & Signal-processing pipeline (5~stages)
  & Deep learning (INR + MMDS loss) \\[6pt]
\textbf{Primary dataset}
  & Chinese Acad.\ Sci.\ (128 subj., 734 pos.)
  & SS2 (78 subj.) \\[6pt]
\textbf{Validation}
  & MUSHRA listening tests (32 dir., 6 LSD levels)
  & Objective perceptual metrics (PBC, AEP, DRMSE) \\[6pt]
\textbf{Key metric}
  & Pearson $r$ with listening-test scores
  & Pearson $r$ between latent \& perceptual distances \\[6pt]
\textbf{LSD critique}
  & No thresholds, no freq.\ weighting, no dir.\ sensitivity
  & Spectral recon.\ error $\neq$ perceptual similarity \\[6pt]
\textbf{Relevance to CPM}
  & Benchmark \& design principles
  & Downstream consumer of composite metrics \\
\bottomrule
\end{tabular}
\end{table}


% =========================================================================
\section{Conclusion}
% =========================================================================

Both papers confirm the central premise of the thesis: LSD (and spectral
reconstruction error more broadly) is an inadequate proxy for perceptual
HRTF quality.  Yao et~al.\ tackle this at the metric level through a
multi-stage auditory pipeline, while Zhang et~al.\ address it at the
representation level by reshaping latent spaces with perceptual
supervision.  Together they bracket the problem space in which the
proposed Composite Perceptual Metric operates, providing both benchmark
comparisons and potential integration pathways.


% =========================================================================
\bibliographystyle{plainnat}
\bibliography{references_report}
% =========================================================================

\end{document}
